\section{从相对论色散到 Dirac 方程}
\label{sec:feqo-dirac-from-dispersion}

为什么自由电子的一章要从 Dirac 方程开始，而不是从薛定谔
方程开始？因为透射电镜的工作电压已经把电子送进相对论区间。
电子静能为 \(mc^2=511\,\si{\kilo\electronvolt}\)；动能
\(K=(\gamma-1)mc^2\) 取 \(200\,\si{\kilo\electronvolt}\) 时
\(\gamma=1+K/(mc^2)=1.391\)，由 \(v=c\sqrt{1-\gamma^{-2}}\)
得 \(v=0.695\,c\)；取 \(300\,\si{\kilo\electronvolt}\) 时
\(\gamma=1.587\)、\(v=0.776\,c\)。动能已达静能的四到六成，
非相对论色散 \(E=mc^2+p^2/(2m)\) 对群速度
\(v=\partial E/\partial p\) 的误差约为 \(v^2/(2c^2)\approx25\%\)，
对后文反复使用的有效质量（色散的二阶导数）误差更大。
动量本身也要用相对论公式：
\(pc=\sqrt{\gamma^2-1}\,mc^2\)，在 \(200\,\si{\kilo\electronvolt}\)
时 \(pc=494\,\si{\kilo\electronvolt}\)，即
\(p_0=\gamma mv=2.64\times10^{-22}\,\si{\kilo\gram\metre\per\second}\)。
所以本章的色散、旋量与电流都是后面各章近场相互作用计算的
运动学底座。

量子力学的时间演化对时间是一阶的，因而初始波函数决定以后各时刻的状态。
相对论能量关系却是 \(E^2=m^2c^4+c^2|\vb{p}|^2\)。
直接把正平方根 \(E(\vb{p})\) 变成微分算子，会得到空间中的非局域算子。
Dirac 的做法是扩大波函数的分量数，用矩阵把这个二次关系分解成
一阶的局域方程。

先写 \(H_0=c\alpha_i p_i+\beta mc^2\)（重复指标 \(i=1,2,3\) 求和），要求
\(H_0^2=(m^2c^4+c^2|\vb{p}|^2)I\) 对每个 \(\vb{p}\)
成立。这一步可以逐条做完。把乘积展开，
\[
H_0^2=c^2\alpha_i\alpha_jp_ip_j
+mc^3(\alpha_i\beta+\beta\alpha_i)p_i+m^2c^4\beta^2 .
\]
因为 \(p_ip_j\) 关于 \(i,j\) 对称，第一项中只有 \(\alpha_i\alpha_j\)
的对称部分有贡献：
\[
\alpha_i\alpha_jp_ip_j
=\frac12(\alpha_i\alpha_j+\alpha_j\alpha_i)p_ip_j .
\]
要让 \(H_0^2\) 对\emph{任意} \(\vb{p}\) 等于
\(c^2|\vb{p}|^2I+m^2c^4I\)，分别比较 \(p_ip_j\)（含 \(i=j\)）、
\(p_i\) 和常数项的系数，得到充要条件
\begin{equation}
\alpha_i\alpha_j+\alpha_j\alpha_i=2\delta_{ij}I,\qquad
\alpha_i\beta+\beta\alpha_i=0,\qquad \beta^2=I.
\label{eq:feqo-dirac-anticommutators}
\end{equation}
这里的反对易子定义为 \(\{A,B\}=AB+BA\)。条件不仅是充分的，
也是必要的：若某一系数不满足，取相应的 \(\vb{p}\)
就能让 \(H_0^2\) 偏离相对论色散。

为什么分量数非得到四？由 \(\alpha_i^2=\beta^2=I\)，
这些矩阵的本征值只能是 \(\pm1\)。又例如
\(\operatorname{Tr}\alpha_i
=\operatorname{Tr}(\beta^2\alpha_i)
=-\operatorname{Tr}(\beta\alpha_i\beta)
=-\operatorname{Tr}\alpha_i\)（第二步用
\(\beta\alpha_i=-\alpha_i\beta\)，最后一步用迹的循环性），
故四个矩阵的迹都为零，维数必须是偶数——否则
\(\pm1\) 本征值无法配成零迹。二维也不行：\(2\times2\)
矩阵里两两反对易的独立矩阵最多只有三个 Pauli 矩阵，
而我们至少需要四个（三个 \(\alpha_i\) 加一个 \(\beta\)）。
所以最小维数是四。

取量子力学中自旋 \(1/2\) 的三个 Pauli 矩阵
\[
\sigma_1=\begin{pmatrix}0&1\\1&0\end{pmatrix},\quad
\sigma_2=\begin{pmatrix}0&-\ii\\\ii&0\end{pmatrix},\quad
\sigma_3=\begin{pmatrix}1&0\\0&-1\end{pmatrix},
\]
它们满足 \(\sigma_i\sigma_j+\sigma_j\sigma_i=2\delta_{ij}I_2\)。
令
\[
\alpha_i=\begin{pmatrix}0&\sigma_i\\\sigma_i&0\end{pmatrix},
\qquad
\beta=\begin{pmatrix}I_2&0\\0&-I_2\end{pmatrix},
\]
逐块相乘便验证\cref{eq:feqo-dirac-anticommutators}。
电子波函数因此是四分量复列向量，满足
\begin{equation}
\ii\hbar\partial_t\Psi
=H_0\Psi,\qquad
H_0=c\boldsymbol\alpha\cdot(-\ii\hbar\nabla)+\beta mc^2.
\label{eq:feqo-dirac-free}
\end{equation}
这并没有把电子的四个分量误当成四个独立粒子；
矩阵关系保证每个平面波仍满足原相对论色散。

这一“\(\boldsymbol\alpha,\beta\) 哈密顿量写法”与量子场论教材的
“\(\gamma^\mu\) 协变写法”是同一个方程，词典只有两行。定义
\begin{equation}
\gamma^0=\beta,\qquad \gamma^i=\beta\alpha_i\quad(i=1,2,3).
\label{eq:feqo-gamma-dictionary}
\end{equation}
用\cref{eq:feqo-dirac-anticommutators} 逐对检验：
\((\gamma^0)^2=\beta^2=I\)；
\((\gamma^i)^2=\beta\alpha_i\beta\alpha_i=-\beta^2\alpha_i^2=-I\)；
\(\gamma^0\gamma^i+\gamma^i\gamma^0
=\alpha_i+\beta\alpha_i\beta=0\)；
\(i\neq j\) 时 \(\gamma^i\gamma^j+\gamma^j\gamma^i
=-\{\alpha_i,\alpha_j\}=0\)。
按全局约定的度规 \(\eta_{\mu\nu}=\operatorname{diag}(1,-1,-1,-1)\)，
这四行合成 Clifford 代数
\begin{equation}
\{\gamma^\mu,\gamma^\nu\}=2\eta^{\mu\nu}I_4 .
\label{eq:feqo-clifford}
\end{equation}
把\cref{eq:feqo-dirac-free} 两边左乘 \(\beta\)，并利用
\(\beta\boldsymbol\alpha\cdot\nabla=\gamma^i\partial_i\)、
\(\beta^2=I\)，就得到协变形式
\begin{equation}
\ii\hbar\gamma^\mu\partial_\mu\Psi=mc\,\Psi,
\qquad \partial_\mu=(\partial_{ct},\nabla).
\label{eq:feqo-dirac-covariant}
\end{equation}
所以凡本书写 \(c\Psi^\dagger\boldsymbol\alpha\Psi\) 的地方，
场论教材写 \(\bar\Psi\gamma^i\Psi\)（\(\bar\Psi=\Psi^\dagger\gamma^0\)）；
两者由\cref{eq:feqo-gamma-dictionary} 直接互译。
这里取的分块形式称为 Dirac（标准）表象；把四个 \(\gamma^\mu\)
同时做相似变换 \(\gamma^\mu\to S\gamma^\mu S^{-1}\) 仍满足
\cref{eq:feqo-clifford}，常用的还有手征（Weyl）表象。
表象只是坐标选择，本节所有物理结果——色散、本征能量、
电流——都不依赖它。


\par
\begin{workedexample}
取 \(\vb{p}=p\hat{\vb{z}}\) 与例题相同的 \(u_1,u_2\)，
用\cref{eq:feqo-gamma-dictionary} 直接计算
\(\bar u_su_{s'}\) 与 \(\sum_su_s\bar u_s\)，
核对正文给出的 \(mc^2/E_{\vb{p}}\delta_{ss'}\) 与
\((c\slashed{p}+mc^2)/(2E_{\vb{p}})\)。
再把归一化换成场论约定 \(u^\dagger u=2E_{\vb{p}}\)
（\(\hbar=c=1\)），写出相应的 \(\bar uu\)。
（提示：换约定只是整体因子；答案为 \(\bar uu=2m\delta_{ss'}\)。）
\textbf{解.} 取 \(\xi_s\) 为 \(\sigma_3\) 本征自旋量，则
\[
u_s=N\begin{pmatrix}\xi_s\\ \frac{cp}{E+mc^2}\sigma_3\xi_s\end{pmatrix},
\qquad N^2=\frac{E+mc^2}{2E}.
\]
于是 \(\bar u_su_{s'}=N^2[1-c^2p^2/(E+mc^2)^2]\delta_{ss'}=(mc^2/E)\delta_{ss'}\)。自旋完备性给 \(\sum_su_s\bar u_s=(c\slashed p+mc^2)/(2E)\)。改用 \(u^\dagger u=2E\)、\(\hbar=c=1\) 时整体放大 \(\sqrt{2E}\)，故 \(\bar uu=2m\delta_{ss'}\)。
\end{workedexample}
\par
\section{正能旋量从本征方程算出}
\label{sec:feqo-positive-spinor}

代入 \(\Psi=u(\vb{p})
\ee^{\ii(\vb{p}\cdot\vb{r}-Et)/\hbar}\)。
写 \(u=(\phi,\chi)^\mathsf T\)，其中上下两块各有两个复分量，
则 \(H_0u=Eu\) 等价于
\[
(E-mc^2)\phi=c\boldsymbol\sigma\cdot\vb{p}\,\chi,\qquad
(E+mc^2)\chi=c\boldsymbol\sigma\cdot\vb{p}\,\phi.
\]
正能根为 \(E_{\vb{p}}=\sqrt{m^2c^4+c^2|\vb{p}|^2}>0\)。
第二式给 \(\chi=c(\boldsymbol\sigma\cdot\vb{p})\phi/
(E_{\vb{p}}+mc^2)\)；代入第一式，利用
\((\boldsymbol\sigma\cdot\vb{p})^2=|\vb{p}|^2I_2\)，
恰还原 \(E_{\vb{p}}^2-m^2c^4=c^2|\vb{p}|^2\)。

任取二分量单位正交基 \(\xi_s\)（\(s=1,2\)），令
\(\phi=N_{\vb{p}}\xi_s\)。归一化条件是
\[
1=u_s^\dagger u_s
=N_{\vb{p}}^2
\left(1+\frac{c^2|\vb{p}|^2}{(E_{\vb{p}}+mc^2)^2}\right)
=N_{\vb{p}}^2\frac{2E_{\vb{p}}}{E_{\vb{p}}+mc^2}.
\]
故
\begin{equation}
u_s(\vb{p})=
\sqrt{\frac{E_{\vb{p}}+mc^2}{2E_{\vb{p}}}}
\begin{pmatrix}
\xi_s\\[1mm]
\dfrac{c\boldsymbol\sigma\cdot\vb{p}}{E_{\vb{p}}+mc^2}\xi_s
\end{pmatrix},
\qquad
u_s^\dagger u_{s'}=\delta_{ss'}.
\label{eq:feqo-positive-spinor}
\end{equation}
对负能根 \(-E_{\vb{p}}\)，先取下分量
\(\chi=N_{\vb{p}}\eta_s\)，第一块方程给
\(\phi=-c(\boldsymbol\sigma\cdot\vb{p})\chi/(E_{\vb{p}}+mc^2)\)。
因此另一组归一化本征旋量为
\[
v_s(\vb{p})=N_{\vb{p}}
\begin{pmatrix}
-\dfrac{c\boldsymbol\sigma\cdot\vb{p}}
{E_{\vb{p}}+mc^2}\eta_s\\[1mm]\eta_s
\end{pmatrix},
\qquad H_0v_s=-E_{\vb{p}}v_s .
\]
其中 \(\eta_s\) 为二分量正交基。
按上下块相乘可查 \(u_s^\dagger v_{s'}=0\)；
正负两组各有两个态，共同张成四维分量空间。
本书多数实验只制备正能电子，但负能子空间在判断投影近似是否有效时不能
从代数上删去。


\par
\begin{workedexample}
由\cref{eq:feqo-positive-spinor} 算
\(u_s^\dagger\boldsymbol\alpha u_s\)，核对
\(\vb{j}/\rho=c^2\vb{p}/E_{\vb{p}}\)。
（提示：上下块相乘后用
\(\sigma_i(\boldsymbol\sigma\cdot\vb{p})
+(\boldsymbol\sigma\cdot\vb{p})\sigma_i=2p_iI_2\)；
归一化因子 \(N_{\vb{p}}\) 在比值中消去。）
\textbf{解.} 令 \(a=cp/(E+mc^2)\)。上下块相乘并用 Pauli 反对易恒等式，得
\[
u_s^\dagger\alpha_i u_s=N_{\vb p}^2\frac{2cp_i}{E+mc^2}
=\frac{c p_i}{E},
\qquad u_s^\dagger u_s=1.
\]
故 Dirac 概率流除以概率密度为 \(c\,u_s^\dagger\boldsymbol\alpha u_s=c^2\vb p/E\)，与色散关系的群速度一致。
\end{workedexample}
\par
归一化约定在这里一次说清，全书不再改动。本书采用
\emph{每粒子概率归一化} \(u_s^\dagger u_{s'}=\delta_{ss'}\)，
使\cref{eq:feqo-positive-wavepacket} 的波包系数直接是概率振幅。
与协变计算配套的量是 Dirac 共轭
\(\bar u=u^\dagger\gamma^0=u^\dagger\beta\)；
用\cref{eq:feqo-positive-spinor} 直接相乘，
\[
\bar u_s(\vb{p})u_{s'}(\vb{p})
=\frac{E_{\vb{p}}+mc^2}{2E_{\vb{p}}}\,
\xi_s^\dagger\!\left(I_2-
\frac{c^2|\vb{p}|^2}{(E_{\vb{p}}+mc^2)^2}\right)\!\xi_{s'}
=\frac{mc^2}{E_{\vb{p}}}\delta_{ss'} .
\]
许多量子场论教材（取 \(\hbar=c=1\)）改用
\(\bar u u=2m\)、\(u^\dagger u=2E_{\vb{p}}\) 的协变归一化，
两种旋量只差因子 \(\sqrt{2E_{\vb{p}}}\)；换算时连同
\(\dd^3p/(2\pi)^3\) 中的因子一起搬即可。

\begin{example}[沿 \(z\) 轴动量的显式旋量]
取 \(\vb{p}=p\hat{\vb{z}}\)，\(\xi_1=(1,0)^\mathsf T\)、
\(\xi_2=(0,1)^\mathsf T\)。此时
\(\boldsymbol\sigma\cdot\vb{p}=p\sigma_3\)，
\(\sigma_3\xi_1=\xi_1\)、\(\sigma_3\xi_2=-\xi_2\)。记
\[
N_{\vb{p}}=\sqrt{\frac{E_{\vb{p}}+mc^2}{2E_{\vb{p}}}},\qquad
r=\frac{cp}{E_{\vb{p}}+mc^2},
\]
则\cref{eq:feqo-positive-spinor} 给出四个分量完全写出的旋量
\[
u_1=N_{\vb{p}}\begin{pmatrix}1\\0\\r\\0\end{pmatrix},\qquad
u_2=N_{\vb{p}}\begin{pmatrix}0\\1\\0\\-r\end{pmatrix}.
\]
检验：\(u_1^\dagger u_1=N_{\vb{p}}^2(1+r^2)\)，而
\(r^2=c^2p^2/(E_{\vb{p}}+mc^2)^2=(E_{\vb{p}}-mc^2)/(E_{\vb{p}}+mc^2)\)
（分子分母同乘 \(E_{\vb{p}}-mc^2\) 并用
\(c^2p^2=E_{\vb{p}}^2-m^2c^4\)），
故 \(1+r^2=2E_{\vb{p}}/(E_{\vb{p}}+mc^2)=N_{\vb{p}}^{-2}\)，
归一化成立；\(u_1^\dagger u_2=0\) 逐分量显然。
取 \(200\,\si{\kilo\electronvolt}\) 电子
（\(\gamma=1.391\)，\(pc=494\,\si{\kilo\electronvolt}\)）：
\(r=pc/(E+mc^2)=494/1022=0.483\)。
下分量并不小——这正是四分量结构在电镜能量下不可省略的定量理由。
\end{example}

每个固定 \(\vb{p}\) 上有两个正能解，这个二重简并就是自旋。
自旋求和可以一次算完：由谱分解
\(\Lambda_+=\sum_su_su_s^\dagger\)（正能投影见下），
\[
\sum_{s=1}^2u_s(\vb{p})\bar u_s(\vb{p})
=\Lambda_+\beta
=\frac{E_{\vb{p}}\gamma^0-c\gamma^ip_i+mc^2I_4}{2E_{\vb{p}}}
=\frac{c\slashed{p}+mc^2I_4}{2E_{\vb{p}}},
\]
其中 \(\slashed{p}=\gamma^\mu p_\mu\)、
\(p^\mu=(E_{\vb{p}}/c,\vb{p})\)。两边取迹并与上式核对：
\(\sum_s\bar u_su_s=2mc^2/E_{\vb{p}}\)，一致。
同理负能侧 \(\sum_sv_s\bar v_s=(c\slashed{p}-mc^2I_4)/(2E_{\vb{p}})\)。
这些自旋和是把矩阵元对未观测自旋求平均时的标准工具。


\begin{proposition}[自由 哈密顿量 的正负能投影]
固定动量 \(\vb{p}\)，定义
\[
\Lambda_\pm(\vb{p})
=\frac12\left(I_4\pm\frac{H_0(\vb{p})}{E_{\vb{p}}}\right).
\]
则 \(\Lambda_\pm^2=\Lambda_\pm\)（幂等）、
\(\Lambda_\pm^\dagger=\Lambda_\pm\)（正交投影）、
\(\Lambda_+\Lambda_-=0\)、\(\Lambda_++\Lambda_-=I_4\)，
\(\operatorname{Tr}\Lambda_\pm=2\)，且
\(H_0\Lambda_\pm=\pm E_{\vb{p}}\Lambda_\pm\)。
\end{proposition}
\begin{proof}
先用\cref{eq:feqo-dirac-anticommutators} 得
\(H_0(\vb{p})^2=E_{\vb{p}}^2I_4\)。
幂等性把平方逐项写出即得：
\[
\left(I\pm\frac{H_0}{E_{\vb{p}}}\right)^2
=I\pm2\frac{H_0}{E_{\vb{p}}}+\frac{H_0^2}{E_{\vb{p}}^2}
=I\pm2\frac{H_0}{E_{\vb{p}}}+I
=2\left(I\pm\frac{H_0}{E_{\vb{p}}}\right),
\]
两边除以 \(4\) 正是 \(\Lambda_\pm^2=\Lambda_\pm\)。
厄米性由 \(\alpha_i^\dagger=\alpha_i\)、\(\beta^\dagger=\beta\)
直接读出。两个不同号的乘积为
\((I+H_0/E_{\vb{p}})(I-H_0/E_{\vb{p}})
=I-H_0^2/E_{\vb{p}}^2=0\)。
迹利用 \(\operatorname{Tr}\alpha_i=\operatorname{Tr}\beta=0\)
（第一节已证），故
\(\operatorname{Tr}\Lambda_\pm=\operatorname{Tr}(I_4)/2=2\)：
每个符号各含两个本征态，对应自旋的二重简并。
最后将 \(H_0\) 乘进定义并再次用平方恒等式。
\end{proof}

投影算符能立即做一件事：对任意四分量动量振幅
\(b(\vb{p})\)，\(\Lambda_+b\) 提取正能部分，
\(\Lambda_-b\) 提取负能部分。若外场随时间变化，两个子空间一般会混合；
只有混合幅度相对能隙足够小，才能把后者作为高阶修正。

负能解为什么不能随手删掉，值得专门说清。在单粒子图像里，
负能态没有下界：一个正能电子若能跃迁到 \(E<0\)
的态并辐射出能量，就会无休止地向下级联，物质将不稳定。
Dirac 的原始补救是“海”的图像：假设所有负能态都被填满，
Pauli 不相容原理挡住向下的跃迁，海中的空穴则表现为
带正电的反粒子（正电子）。这个图像在逻辑上不自洽——
它无法解释为什么无穷大的背景电荷和能量不可观测，
更根本的是，\(\hbar\omega>2mc^2=1.022\,\si{\mega\electronvolt}\)
的场可以直接从真空中产生正负电子对，此时“粒子数守恒的
单粒子方程”这一前提本身失效。现代的处理是把 Dirac
方程重新解释为经典场方程，再做二次量子化：
负能解的展开系数变成正电子的产生算符，负能困难转化为
反粒子的存在。这条路的系统展开见标准量子场论教材
（如 Peskin \& Schroeder 第 1--3 章）。本书不需要走那么远：
光频光子的能量约 \(1.5\,\si{\electronvolt}\)（
\(800\,\si{\nano\metre}\)），与对产生阈
\(1.022\,\si{\mega\electronvolt}\) 差六个数量级；
电镜中电子密度也极低，一次只有一个电子参与相互作用。
因此全书在正能子空间内工作，\(\Lambda_-\)
只用于检验这个投影近似的有效性（见\cref{sec:feqo-minimal-coupling}
末尾的定量估计）。

四维分量里有两个是自旋，这一点后文如何使用也要交代。
自旋自由度通过 \(\boldsymbol\alpha\) 与动量耦合，
凡涉及磁场、自旋轨道耦合或大角弹性散射（Mott 散射）
的过程，两个自旋态的演化不同，必须保留完整的旋量结构。
本书处理的却是另一类相互作用：纳米结构的近场在电子
德布罗意波长（\(200\,\si{\kilo\electronvolt}\) 时约
\(2.5\,\si{\pico\metre}\)）的尺度上近似均匀，相互作用
矩阵元在两个自旋态之间近似对角且相等，自旋因此成为一个
整体的二重简并标签。从\cref{chap:feqo-finite-time} 起，
各章都固定一个自旋态、把波函数写成标量包络；
需要自旋求和时，用上面的
\(\sum_su_s\bar u_s=(c\slashed{p}+mc^2)/(2E_{\vb{p}})\)
即可恢复完整结构。


\section{概率守恒与波包归一化}
\label{sec:feqo-current-and-packet}

对\cref{eq:feqo-dirac-free} 左乘 \(\Psi^\dagger\)，
再取其厄米共轭并右乘 \(\Psi\)，质量项因 \(\beta^\dagger=\beta\)
相消；\(\alpha_i^\dagger=\alpha_i\) 使空间导数合成散度：
\[
\partial_t(\Psi^\dagger\Psi)
=-c\bigl[(\nabla\Psi^\dagger)\cdot\boldsymbol\alpha\Psi
+\Psi^\dagger\boldsymbol\alpha\cdot\nabla\Psi\bigr]
=-\nabla\cdot(c\Psi^\dagger\boldsymbol\alpha\Psi).
\]
因此定义
\begin{equation}
\rho=\Psi^\dagger\Psi,\qquad
\vb{j}=c\Psi^\dagger\boldsymbol\alpha\Psi,\qquad
\partial_t\rho+\nabla\cdot\vb{j}=0.
\label{eq:feqo-dirac-current}
\end{equation}
\(\rho\) 的单位是 \(\si{\per\cubic\metre}\)，
\(\vb{j}\) 是概率流而非电荷流；电子电荷流为 \(q\vb{j}\)。
这个区分有实际后果：进入 Maxwell 方程、与矢势耦合的是
电荷流 \(q\vb{j}\)（见\cref{sec:feqo-minimal-coupling}），
而归一化条件约束的是概率流；在非相对论极限下，
\(\vb{j}\) 退化为薛定谔理论里熟悉的
\(\hbar\operatorname{Im}(\Psi^*\nabla\Psi)/m\)。

用\cref{eq:feqo-gamma-dictionary} 的词典，
\cref{eq:feqo-dirac-current} 正是场论中守恒四矢量的分量写法。
定义
\[
j^\mu=c\bar\Psi\gamma^\mu\Psi,\qquad
\bar\Psi=\Psi^\dagger\gamma^0 .
\]
时间分量 \(j^0=c\Psi^\dagger\gamma^0\gamma^0\Psi=c\rho\)；
空间分量 \(j^i=c\Psi^\dagger\gamma^0\gamma^i\Psi
=c\Psi^\dagger\alpha_i\Psi\)，即 \(\vb{j}\)。
所以 \(j^\mu=(c\rho,\vb{j})\)，连续性方程紧凑地写成
\(\partial_\mu j^\mu=0\)。这正是量子电动力学中
电子—光子顶角所用的同一个四矢量；
单粒子层面它已出现在跃迁矩阵元中
（\cref{sec:feqo-minimal-coupling} 的跃迁电流）。

若波函数在无穷远衰减，使边界通量为零，对全空间积分便得到
\(\dd\int\rho\,\dd^3r/\dd t=0\)。

对正能平面波，还能核对流速。把\cref{eq:feqo-positive-spinor}
代入 \(u_s^\dagger\alpha_i u_s\)，上下块相乘后利用
\(\sigma_i(\boldsymbol\sigma\cdot\vb{p})
+(\boldsymbol\sigma\cdot\vb{p})\sigma_i=2p_iI_2\)，得到
\[
c\,u_s^\dagger\alpha_i u_s
=\frac{c^2p_i}{E_{\vb{p}}}
=\frac{\partial E_{\vb{p}}}{\partial p_i}.
\]
概率流除以概率密度，正是色散的群速度。

平面波不能在全空间归一化。全书取
\[
\langle\vb{r}|\vb{p},s\rangle
=(2\pi\hbar)^{-3/2}u_s(\vb{p})
\ee^{\ii\vb{p}\cdot\vb{r}/\hbar}.
\]
用 Fourier 恒等式
\(\int\dd^3r\,\ee^{\ii(\vb{p}'-\vb{p})\cdot\vb{r}/\hbar}
=(2\pi\hbar)^3\delta^{(3)}(\vb{p}'-\vb{p})\)，
并在 delta 支持处使用 \(u_s^\dagger u_{s'}=\delta_{ss'}\)，得
\(\langle\vb{p},s|\vb{p}',s'\rangle
=\delta_{ss'}\delta^{(3)}(\vb{p}-\vb{p}')\)。
因此归一化的正能波包为
\begin{equation}
\Psi(\vb{r},t)=
\sum_s\int\frac{\dd^3p}{(2\pi\hbar)^{3/2}}
a_s(\vb{p})u_s(\vb{p})
\ee^{\ii(\vb{p}\cdot\vb{r}-E_{\vb{p}}t)/\hbar},
\qquad
\sum_s\int\dd^3p\,|a_s(\vb{p})|^2=1.
\label{eq:feqo-positive-wavepacket}
\end{equation}
初始动量宽度、能量宽度和脉冲长度由 \(a_s\) 决定，
不能在传播或探测阶段重新当作自由参数调节。


\par
\begin{workedexample}
设波包主要满足 \(|\delta p_z|\le\Delta p\)。
用本节的余项界写出传播时间 \(t\) 后二阶近似相位误差小于
\(0.1\) 的充分条件，并说明缩短 \(t\) 与缩小 \(\Delta p\)
各会带来什么改变。
（提示：充分条件形如 \(tM_3(\Delta p)^3/(6\hbar)<0.1\)；
注意 \(\Delta p\) 以三次方进入，压缩带宽比缩短时间有效得多。）
\textbf{解.} Taylor 余项满足 \(|R_3|\le M_3(\Delta p)^3/6\)，传播相位误差不超过 \(t|R_3|/\hbar\)。因此充分条件为
\[
\frac{tM_3(\Delta p)^3}{6\hbar}<0.1.
\]
将 \(t\) 减半使误差界减半；将 \(\Delta p\) 减半则使误差界降为原来的 \(1/8\)，同时会使波包更长。
\end{workedexample}
\par
\section{窄带传播究竟舍弃了什么}
\label{sec:feqo-dispersion-expansion}

取中心动量 \(\vb{p}_0=p_0\hat{\vb{z}}\)，写
\(\delta\vb{p}=\vb{p}-\vb{p}_0\)、
\(E_0=E_{\vb{p}_0}\)、
\(\gamma_0=E_0/(mc^2)\)。
对精确色散求导，
\[
\partial_iE=\frac{c^2p_i}{E},\qquad
\partial_i\partial_jE
=\frac{c^2}{E}\delta_{ij}
-\frac{c^4p_ip_j}{E^3}.
\]
在 \(\vb{p}_0\) 处，纵向二阶导数为
\(1/(\gamma_0^3m)\)，两个横向二阶导数为
\(1/(\gamma_0m)\)，混合导数为零。因此
\begin{equation}
E(\vb{p})
=E_0+v_0\delta p_z
+\frac{\delta p_\perp^2}{2\gamma_0m}
+\frac{\delta p_z^2}{2\gamma_0^3m}
+R_3(\delta\vb{p}).
\label{eq:feqo-dispersion-hessian}
\end{equation}

纵向与横向的有效质量并不相等。
\cref{fig:relativistic-effective-mass} 把两条质量比曲线画在同一张图上，
横轴是动能 \(K\)，纵轴是质量比：\(\gamma\) 是横向的 \(m_\perp/m_e\)，
\(\gamma^3\) 是纵向的 \(m_\parallel/m_e\)。
两条曲线都从 \(K=0\) 处的 \(1\) 出发，但 \(\gamma^3\) 上升得远快得多；
在 \(K=200\,\si{\kilo\electronvolt}\) 处 \(\gamma\approx1.39\)、
\(\gamma^3\approx2.69\)，纵向有效质量已接近横向的两倍。
这就是\cref{eq:feqo-dispersion-hessian} 中
\(\delta p_z^2\) 与 \(\delta p_\perp^2\) 的系数
\(1/(\gamma_0^3m)\) 与 \(1/(\gamma_0m)\) 相差一个 \(\gamma_0^2\) 的图形来源。

\begin{figure}[htbp]
\centering
\includegraphics[width=0.76\linewidth]{relativistic_effective_mass.pdf}
\caption{纵向与横向有效质量的能量依赖。横轴为动能 $K$，纵轴为质量比：$\gamma$ 是横向质量比 $m_\perp/m_e$，$\gamma^3$ 是纵向质量比 $m_\parallel/m_e$。两者在 $K=0$ 处同为 $1$，随能量升高 $\gamma^3$ 增长明显更快。}
\label{fig:relativistic-effective-mass}
\end{figure}

若波包支集位于中心周围半径 \(\Delta p\) 的球内，
且该球中三阶导数的算子范数不超过 \(M_3\)，Taylor 定理给
\(|R_3|\le M_3(\Delta p)^3/6\)。
传播时间为 \(t\) 时，由舍去余项造成的相位误差至多
\(tM_3(\Delta p)^3/(6\hbar)\)。
这个无量纲数远小于 \(1\) 时二阶包络才可靠；
“高能电子”本身不是省略三阶项的理由。

这个余项界在具体能量上到底有多大，可以直接看图。
\cref{fig:dispersion-control-modern} 取 200\,\si{\kilo\electronvolt} 电子，
横轴是相对动量偏离 \(\delta p/p_0\in[-0.38,0.38]\)，
纵轴是 \((E-E_0)/(m_ec^2)\)，把精确相对论色散、一阶展开
\(v_0\delta p\) 与二阶展开 \(v_0\delta p+\delta p^2/(2m_\parallel)\)
画在同一张图上。
一阶直线在 \(|\delta p/p_0|\approx0.2\) 处就与精确曲线明显分开，
二阶抛物线则几乎在整个窗口内贴着精确曲线，
只在窗口边缘留下约纵轴量程百分之一的残差；又因三阶项是奇函数，
这份残差在窗口两侧符号相反。

\begin{figure}[htbp]
\centering
\includegraphics[width=0.78\linewidth]{dispersion_control.pdf}
\caption{200\,\si{\kilo\electronvolt} 电子的精确相对论色散与在 $p_0$ 处的一阶、二阶展开。横轴为 $\delta p/p_0$，纵轴为 $(E-E_0)/(m_ec^2)$；二阶曲线偏离精确曲线的位置就是高阶色散开始起作用的地方。}
\label{fig:dispersion-control-modern}
\end{figure}

二阶近似能保留多少传播信息？用一维高斯电子包从头算一次。
固定一个正能自旋态并只看纵向。以 \(p_0\) 为中心的归一化
动量包络取为
\[
a(p)=(2\pi\sigma_p^2)^{-1/4}
\exp\!\left[-\frac{(p-p_0)^2}{4\sigma_p^2}\right],
\qquad \int_{-\infty}^{\infty}|a(p)|^2\,\dd p=1.
\]
假设 \(\sigma_p\ll p_0\)，负动量尾部才可忽略，
这使单向传播的读法与积分范围相容。
“\(\ll\)”可以给出具体数字。能量宽度与动量宽度由
\(\sigma_E=v_0\sigma_p\) 联系（色散的一阶导数）。
取 \(200\,\si{\kilo\electronvolt}\) 电子：
\(p_0=2.64\times10^{-22}\,\si{\kilo\gram\metre\per\second}\)、
\(v_0=0.695\,c\)。电镜束流的典型能量宽度
\(\sigma_E\sim1\,\si{\electronvolt}\) 给
\(\sigma_p=\sigma_E/v_0\approx7.7\times10^{-28}\,
\si{\kilo\gram\metre\per\second}\)，即
\(\sigma_p/p_0\approx3\times10^{-6}\)；
即使把束流恶化到 \(\sigma_E/E=10^{-4}\)，仍有
\(\sigma_p/p_0\approx6\times10^{-5}\)。
负动量尾部 \(e^{-p_0^2/(4\sigma_p^2)}\) 在两种情形下都
小到无法测量，窄带假设对电镜束流极为宽松；
真正的限制来自前面 Taylor 余项界中的传播时间。
设 \(m_\parallel=\gamma_0^3m\)，
并在\cref{eq:feqo-dispersion-hessian} 中只保留
\(E(p)=E_0+v_0\delta p+\delta p^2/(2m_\parallel)\)。
同时把随动量缓慢变化的旋量 \(u_s(p)\) 取为 \(u_s(p_0)\)；
在有效动量窗口中，这一步的幅度误差受
\(\sigma_p\sup\|\partial_pu_s\|\) 控制。
把此式与 \(a(p)\) 代入一维 Fourier 积分，令
\(\xi=z-v_0t\) 与 \(\sigma_{z0}=\hbar/(2\sigma_p)\)，
待积分的指数就是
\[
-\frac{\delta p^2}{4\sigma_p^2}
\left(1+\ii\frac{2\sigma_p^2t}{m_\parallel\hbar}\right)
+\frac{\ii\delta p\xi}{\hbar}.
\]
对 \(\delta p\) 配方并使用高斯积分，得到略去整体载波
\(\ee^{\ii(p_0z-E_0t)/\hbar}\) 后的包络
\[
\psi_{\rm env}(\xi,t)
=\frac{1}{(2\pi\sigma_{z0}^2)^{1/4}}
\frac{1}{\sqrt{1+\ii\tau}}
\exp\!\left[-\frac{\xi^2}{4\sigma_{z0}^2(1+\ii\tau)}\right],
\qquad
\tau=\frac{\hbar t}{2m_\parallel\sigma_{z0}^2}.
\]
取模平方时用
\(\operatorname{Re}(1+\ii\tau)^{-1}=(1+\tau^2)^{-1}\)，
得到可直接积分检验归一化的密度
\[
|\psi_{\rm env}(\xi,t)|^2
=\frac{\exp[-\xi^2/(2\sigma_{z0}^2(1+\tau^2))]}
{\sqrt{2\pi}\,\sigma_{z0}\sqrt{1+\tau^2}}.
\]
所以位置方差为
\[
\sigma_z^2(t)=\sigma_{z0}^2(1+\tau^2)
=\sigma_{z0}^2+
\frac{\hbar^2t^2}{4m_\parallel^2\sigma_{z0}^2}.
\]
中心按群速度移动，宽度则由纵向有效质量 \(\gamma_0^3m\)
控制。代入数字可以给这个公式一把尺子。宽度增长一倍
（\(\tau=1\)）所需的时间是
\(t_2=2m_\parallel\sigma_{z0}^2/\hbar\)。
取 \(200\,\si{\kilo\electronvolt}\)
（\(m_\parallel=\gamma_0^3m=2.69\,m\)）：
若波包初宽 \(\sigma_{z0}=1\,\si{\nano\metre}\)，则
\(t_2\approx4.7\times10^{-14}\,\si{\second}\)，
其间电子飞行 \(v_0t_2\approx9.7\,\si{\micro\metre}\)；
若 \(\sigma_{z0}=10\,\si{\nano\metre}\)，时间放大一百倍到
\(t_2\approx4.7\,\si{\pico\second}\)，飞行距离约
\(1\,\si{\milli\metre}\)。
纵向宽度平方反比地决定展宽速率，而展宽的“时钟”由
\(\gamma_0^3\) 而非 \(\gamma_0\) 控制——用非相对论质量
会把这个时间低估到三分之一。后文讨论阿秒聚束时，
正是这个展宽时间决定了自由传播段能允许多长。这个结果依赖窄带二阶近似和固定自旋；
高斯分布虽没有严格有限的支集，但在
\(|\delta p|\le L\sigma_p\) 内可用前述 Taylor 余项界，
外部概率由高斯尾积分直接控制。
若 \(tM_3(L\sigma_p)^3/(6\hbar)\) 不小，
或选定 \(L\) 后的尾概率不可忽略，
就应回到精确色散的积分，不能把上式外推到任意远的传播距离。

\section{电磁势怎样进入 哈密顿量}
\label{sec:feqo-minimal-coupling}

电子电荷写作 \(q=-e\)，其中 \(e>0\)。
经典带电粒子的拉格朗日量是
\[
L=-mc^2\sqrt{1-v^2/c^2}
+q\vb{v}\cdot\vb{A}(\vb{r},t)-q\phi(\vb{r},t).
\]
定义 Lorentz 因子
\(\gamma=(1-v^2/c^2)^{-1/2}\)。
由 \(\vb{p}=\partial L/\partial\vb{v}\) 得
\(\vb{p}=\gamma m\vb{v}+q\vb{A}\)：
正则动量与机械动量 \(\gamma m\vb{v}\) 相差 \(q\vb{A}\)。
Legendre 变换给
\[
H_{\rm cl}
=\sqrt{m^2c^4+c^2|\vb{p}-q\vb{A}|^2}+q\phi.
\]
把同一个机械动量替换用于一阶 Dirac 哈密顿量，得到
\begin{equation}
H_D=c\boldsymbol\alpha\cdot
(-\ii\hbar\nabla-q\vb{A})+\beta mc^2+q\phi .
\label{eq:feqo-minimal-coupling}
\end{equation}
这也固定了符号：电子 \(q=-e\)，因此不能把
\(\vb{p}-q\vb{A}\) 随手写成 \(\vb{p}-e\vb{A}\)。

规范变换
\(\vb{A}'=\vb{A}+\nabla\chi\)、
\(\phi'=\phi-\partial_t\chi\) 后，令
\(\Psi'=\ee^{\ii q\chi/\hbar}\Psi\)。
直接对指数求导便有
\[
(-\ii\hbar\nabla-q\vb{A}')\Psi'
=\ee^{\ii q\chi/\hbar}
(-\ii\hbar\nabla-q\vb{A})\Psi,
\]
时间导数与 \(q\phi'\) 中的 \(\partial_t\chi\) 同样相消。
所以同一概率流和实验预测不会随势的规范表示改变。

上面的连续性方程只写了同一波函数的概率守恒。
以后求跃迁振幅时出现的是两个不同电子态之间的电流，
因而还需要它的双线性版本。设 \(\Psi_i,\Psi_f\)
都是同一实电磁势下 Dirac 方程的解，定义
\[
\rho_{fi}^{\rm ch}=q\Psi_f^\dagger\Psi_i,\qquad
\vb{j}_{fi}^{\rm ch}=qc\Psi_f^\dagger
\boldsymbol\alpha\Psi_i .
\]
这两个量一般是复数，并不是一次实验中的实电荷密度或实电流；
它们是跃迁矩阵元。把 \(\Psi_i\) 的方程左乘
\(\Psi_f^\dagger\)，再从 \(\Psi_f\) 方程的共轭式减去，
同一标势与矢势项逐项抵消，得到
\begin{equation}
\partial_t\rho_{fi}^{\rm ch}
+\nabla\cdot\vb{j}_{fi}^{\rm ch}=0 .
\label{eq:feqo-transition-current-continuity}
\end{equation}
若两态是自由定态，时间因子为
\(\exp[-\ii(E_i-E_f)t/\hbar]\)。
记跃迁频率 \(\omega_{if}=(E_i-E_f)/\hbar\)，
则\cref{eq:feqo-transition-current-continuity}
在这一频率分量上读作
\(\nabla\cdot\vb{j}_{fi}^{\rm ch}
=\ii\omega_{if}\rho_{fi}^{\rm ch}\)。
所以纵向电流和电荷密度不能在跃迁计算中独立指定。

这条关系也给规范不变性一个直接的检验。
对一个守恒的跃迁源，势的线性耦合为
\(H_{fi}(t)=\int(\rho_{fi}^{\rm ch}\phi
-\vb{j}_{fi}^{\rm ch}\cdot\vb{A})\,\dd^3r\)。
规范变换使它改变
\[
\delta H_{fi}
=-\int\!\bigl(\rho_{fi}^{\rm ch}\partial_t\chi
+\vb{j}_{fi}^{\rm ch}\cdot\nabla\chi\bigr)\,\dd^3r
=-\frac{\dd}{\dd t}
\int\rho_{fi}^{\rm ch}\chi\,\dd^3r ,
\]
最后一步用\cref{eq:feqo-transition-current-continuity}
并假定空间边界项消失。总时间导数只改变端点态的相位，
不会改变完整过程的跃迁概率。如果只保留
\(-\vb{j}\cdot\vb{A}\) 却在有纵向场时扔掉
\(\rho\phi\)，这个检验会失败；后文使用纯辐射规范时
必须说明为何可以这样选取。

规范自由在近场问题里有一个实用读法。纳米结构的准静态近场
由电荷分布产生，最自然的写法是纯标势描述
（\(\vb{A}=0\)，耦合全在 \(q\phi\) 里）；
也可以做一次规范变换把 \(\phi\) 完全消去，
代价是引入含时的 \(\vb{A}\)，耦合改由
\(q\vb{v}\cdot\vb{A}\) 承担。两种写法给出相同的相位
\((q/\hbar)\int(\vb{v}\cdot\vb{A}-\phi)\,\dd t\)，
因为被积函数只差一个全微分。比较不同文献的耦合公式时，
先确认对方选了哪种规范，再比较数值。

这里的双线性电流还有一个值得知道的身份。在量子电动力学中，
电子—光子相互作用的 Feynman 顶角正是
\(\bar u(\vb{p}')\gamma^\mu u(\vb{p})\) 乘光子极化矢量；
用\cref{eq:feqo-gamma-dictionary} 翻译，其空间分量就是
\(u^\dagger(\vb{p}')\boldsymbol\alpha u(\vb{p})\)。
本节的跃迁电流是这个顶角的单粒子版本；
只要正负电子对产生可以忽略（条件见下），两者的物理内容相同。

外场若含有足够强的频率或空间变化，就可能把
\(\Lambda_+\) 子空间耦合到 \(\Lambda_-\)；
把单电子正能动力学当作封闭模型需要比较
\(\|\Lambda_-H_I\Lambda_+\|\) 与最小正负能隙，
并检查外场频谱是否跨越该能隙。强场对产生区域必须使用量子场描述，
不能把\cref{eq:feqo-minimal-coupling} 的单粒子解释无限延伸。

这个条件可以当场算出数量级。混合矩阵元的尺度是外场在
Compton 波长 \(\bar{\lambda}_C=\hbar/(mc)=3.86\times10^{-13}\,
\si{\metre}\) 上做的功：
\(\|\Lambda_-H_I\Lambda_+\|\sim eE_*\bar{\lambda}_C\)，
而能隙是 \(2mc^2=1.022\,\si{\mega\electronvolt}\)。
取 PINEM 实验中已算强的近场振幅
\(E_*=10^{7}\,\si{\volt\per\metre}\)：
\(eE_*\bar{\lambda}_C\approx3.9\,\si{\micro\electronvolt}\)，
与能隙之比约 \(4\times10^{-12}\)。
即使取原子尺度的 \(10^{11}\,\si{\volt\per\metre}\)，
比值也只有 \(4\times10^{-8}\)。
频谱一侧同样安全：驱动频率
\(\hbar\omega\approx1.5\,\si{\electronvolt}\)
距对产生阈 \(2mc^2\) 六个数量级。
正能投影在本书全部应用场景中都是极佳近似；
它的失效恰恰标志着必须改用电子场的二次量子化描述。



